\documentclass[11pt,a4paper]{article}

\usepackage[T1]{fontenc}
\usepackage[utf8]{inputenc}
\usepackage{lmodern}
\usepackage{amssymb,amsmath,amsthm,mathtools}
\usepackage[margin=2.5cm]{geometry}
\usepackage{parskip}
\usepackage[hidelinks]{hyperref}
\usepackage{microtype}

% Theorem environments (arXiv / journal style)
\theoremstyle{plain}
\newtheorem{theorem}{Theorem}[section]
\newtheorem{proposition}[theorem]{Proposition}
\newtheorem{lemma}[theorem]{Lemma}
\newtheorem{corollary}[theorem]{Corollary}

\theoremstyle{definition}
\newtheorem{definition}[theorem]{Definition}
\newtheorem{example}[theorem]{Example}
\newtheorem{remark}[theorem]{Remark}

% Notation
\newcommand{\Type}{\mathbf{Type}}
\newcommand{\tauu}{\tau}

\title{Composing Trajectory Evaluation with Trust Evaluation}
\author{Extracted from a machine-checked development}
\date{}

\begin{document}
\maketitle

\begin{abstract}
Trajectory systems and trust systems grade claims in different chains. We
define a seven-level trust chain and a monotone map $\psi$ from the
five-level trajectory chain into it, sending bottoms to bottoms. Composing
$\psi$ with canonical trust on the circular and honest trajectories yields
concrete transport results: circular evaluation maps entirely to
$\mathsf{unchecked}$, while honest evaluation maps the verify event to
$\mathsf{proved}$ and the claim event to $\mathsf{computed}$. The proofs of
all stated results have been machine-checked; we record only the mathematical
content.
\end{abstract}

\tableofcontents
\bigskip

%======================================================================
\section{Introduction}
%======================================================================

A trajectory system grades agent actions in a five-level chain specialised
to verification workflows. A trust system grades mathematical claims in a
(typically longer) chain of check strengths used by research agents. To
compose the two evaluations---to read a trajectory grade as a trust
grade---one needs a map of chains that preserves order and bottom.

In the general theory of trust systems, morphisms transport canonical trust
laxly: if $\phi$ is a monotone map of chains with $\phi(\bot) = \bot$, and
$F$ maps claims and steps coherently, then
\begin{equation}
  \phi\bigl(\tauu(x)\bigr) \;\le\; \tauu'\bigl(F(x)\bigr).
\end{equation}
The present note does not re-prove that abstract transport lemma. Instead it
exhibits a concrete chain map $\psi$ from trajectory levels to trust levels
and evaluates $\psi \circ \tauu$ on the two running examples of the companion
note on trajectories: the circular trajectory (where $\tauu \equiv \bot$)
and the honest trajectory (where $\tauu$ attains independent grades along
sound paths).

The material is elementary. Proofs are indicated by the phrase ``by lean,''
meaning they have been checked by a proof assistant rather than written out
in prose.

\paragraph{What this note claims.}
In ordinary mathematical language, the claim is: there is a bottom-preserving
monotone map $\psi$ from the trajectory chain to the seven-level trust chain
such that
\begin{itemize}
  \item every event of the circular trajectory evaluates, after $\psi$, to
        $\mathsf{unchecked}$;
  \item the honest verify event evaluates to $\mathsf{proved}$, and the
        honest claim event to $\mathsf{computed}$.
\end{itemize}
Trajectory evaluation therefore composes with trust evaluation by pure
postcomposition of grades; the numerical content of transport is completely
explicit on the examples.

%======================================================================
\section{Recollections}
%======================================================================

We recall the trajectory chain and the two systems from the companion note;
full definitions of trust systems and canonical trust appear there and in
the note on trust systems.

\begin{definition}[Trajectory chain]
The \textbf{trajectory chain} $T_{\mathrm{traj}}$ is
\begin{equation}
  \mathsf{unsupported}
  \;<\;
  \mathsf{shallow}
  \;<\;
  \mathsf{independent}
  \;<\;
  \mathsf{crossVerified}
  \;<\;
  \mathsf{mechChecked},
\end{equation}
with bottom $\mathsf{unsupported}$.
\end{definition}

\begin{definition}[Events and systems]
Let $E = \{e_0, e_1, e_2, e_3, e_4\}$. The \textbf{circular trajectory}
$\mathcal{S}_{\mathrm{circ}}$ and the \textbf{honest trajectory}
$\mathcal{S}_{\mathrm{hon}}$ are trust systems on $E$ graded in
$T_{\mathrm{traj}}$, sharing the dependency topology
\begin{equation}
  e_0 \to e_1 \to e_2 \to e_3,
  \qquad
  e_2 \to e_4,
\end{equation}
and differing only in grades. Writing $\tauu_{\mathrm{circ}}$ and
$\tauu_{\mathrm{hon}}$ for their canonical trusts, the companion note
establishes
\begin{align}
  \tauu_{\mathrm{circ}}(e) &\;=\; \mathsf{unsupported}
  \quad\text{for all } e \in E, \\
  \tauu_{\mathrm{hon}}(e_3) &\;=\; \mathsf{shallow}, \\
  \tauu_{\mathrm{hon}}(e_4) &\;=\; \mathsf{independent}.
\end{align}
(Along the honest spine $e_0, e_1, e_2$ one likewise has
$\tauu_{\mathrm{hon}} = \mathsf{independent}$.)
\end{definition}

%======================================================================
\section{The trust chain}
%======================================================================

\begin{definition}[Trust levels]\label{def:trustlevel}
The \textbf{trust chain} $T_{\mathrm{trust}}$ is the seven-element chain
\begin{equation}
\begin{aligned}
  &\mathsf{unchecked}
  \;<\;
  \mathsf{speculative}
  \;<\;
  \mathsf{computed}
  \;<\;
  \mathsf{proved} \\
  &\qquad
  \;<\;
  \mathsf{peerReviewed}
  \;<\;
  \mathsf{published}
  \;<\;
  \mathsf{leanVerified},
\end{aligned}
\end{equation}
with bottom $\bot = \mathsf{unchecked}$.
\end{definition}

\begin{proposition}[Chain structure]\label{prop:trust-chain}
$T_{\mathrm{trust}}$ carries the structure of a bounded chain: $\le$ is a
linear order, $\bot$ is least, and $\le$ is decidable.
\end{proposition}

\begin{proof}
By lean.
\end{proof}

\begin{remark}
Informally these levels grade mathematical claims in a research-agent
deployment: from no check, through speculative and computed evidence, to
proved, peer-reviewed, published, and machine-verified claims. The precise
operational meaning of each level is not needed for the theorems below;
only the order and the bottom element matter.
\end{remark}

%======================================================================
\section{The trajectory--trust map}
%======================================================================

\begin{definition}[The map $\psi$]\label{def:psi}
Define $\psi : T_{\mathrm{traj}} \to T_{\mathrm{trust}}$ by
\begin{center}
\begin{tabular}{ll}
\hline
Trajectory level & Trust level \\
\hline
$\mathsf{unsupported}$ & $\mathsf{unchecked}$ \\
$\mathsf{shallow}$ & $\mathsf{computed}$ \\
$\mathsf{independent}$ & $\mathsf{proved}$ \\
$\mathsf{crossVerified}$ & $\mathsf{peerReviewed}$ \\
$\mathsf{mechChecked}$ & $\mathsf{leanVerified}$ \\
\hline
\end{tabular}
\end{center}
\end{definition}

\begin{remark}
The map skips $\mathsf{speculative}$ and $\mathsf{published}$: trajectory
levels are coarser than trust levels, and $\psi$ lands on the trust grades
that best match the intended strength of each trajectory grade. In
particular, independent trajectory checks are read as proved claims, and
mechanically checked trajectory steps as lean-verified claims.
\end{remark}

\begin{theorem}[Bottom preservation]\label{thm:psi-bot}
\begin{equation}
  \psi(\bot_{T_{\mathrm{traj}}}) \;=\; \bot_{T_{\mathrm{trust}}}.
\end{equation}
Equivalently, $\psi(\mathsf{unsupported}) = \mathsf{unchecked}$.
\end{theorem}

\begin{proof}
By lean.
\end{proof}

\begin{theorem}[Monotonicity]\label{thm:psi-mono}
For all $a, b \in T_{\mathrm{traj}}$,
\begin{equation}
  a \le b \;\Longrightarrow\; \psi(a) \le \psi(b).
\end{equation}
\end{theorem}

\begin{proof}
By lean.
\end{proof}

\begin{remark}[Chain-morphism requirements]
In the abstract theory of morphisms of trust systems, a map of check-posets
must be monotone and strict at bottom. Theorems~\ref{thm:psi-bot}
and~\ref{thm:psi-mono} record exactly those two requirements for $\psi$.
Together they make $\psi$ eligible as the grade component of a trust-system
morphism from a trajectory system into a trust system graded in
$T_{\mathrm{trust}}$.
\end{remark}

%======================================================================
\section{Concrete transport}
%======================================================================

Composing $\psi$ with canonical trust on the two examples yields fully
evaluated trust grades---no inequality, only equality.

\begin{theorem}[Circular trajectory transports to unchecked]
\label{thm:psi-circular}
For every event $e \in E$,
\begin{equation}
  \psi\bigl(\tauu_{\mathrm{circ}}(e)\bigr)
  \;=\;
  \mathsf{unchecked}.
\end{equation}
\end{theorem}

\begin{proof}
By lean.
\end{proof}

\begin{remark}
This is immediate from circular collapse
($\tauu_{\mathrm{circ}} \equiv \mathsf{unsupported}$) and bottom
preservation of $\psi$. It is recorded as a theorem because it is the
operational statement used by deployments: after transport, a circular
trajectory contributes no trust above unchecked at any event.
\end{remark}

\begin{theorem}[Honest verify event]\label{thm:psi-honest-e4}
\begin{equation}
  \psi\bigl(\tauu_{\mathrm{hon}}(e_4)\bigr)
  \;=\;
  \mathsf{proved}.
\end{equation}
\end{theorem}

\begin{proof}
By lean.
\end{proof}

\begin{theorem}[Honest claim event]\label{thm:psi-honest-e3}
\begin{equation}
  \psi\bigl(\tauu_{\mathrm{hon}}(e_3)\bigr)
  \;=\;
  \mathsf{computed}.
\end{equation}
\end{theorem}

\begin{proof}
By lean.
\end{proof}

\begin{remark}[Reading the honest results]
On the honest trajectory, $\tauu(e_4) = \mathsf{independent}$ and
$\tauu(e_3) = \mathsf{shallow}$. Applying $\psi$ sends independent to proved
and shallow to computed. Thus a sound verify branch is trusted at the proved
level after composition, while a shallow claim branch is trusted only at the
computed level---matching the intended gap between verification and
assertion.
\end{remark}

%======================================================================
\section{Discussion}
%======================================================================

\paragraph{What has been established.}
Trajectory grades embed into trust grades via a bottom-preserving monotone
map $\psi$. On the circular example the composite is constantly unchecked;
on the honest example the verify and claim events map to proved and
computed respectively. Composition of the two evaluations is therefore a
matter of postcomposing the grade map---no further analysis of dependency
structure is required once $\tauu$ is known.

\paragraph{Relation to abstract transport.}
The general transport lemma for morphisms of trust systems asserts a
\emph{lax} inequality $\phi(\tauu(x)) \le \tauu'(F(x))$. The results here
are stronger on the examples because they evaluate $\psi \circ \tauu$
exactly, without constructing a target system $S'$ or a claim map $F$. They
may be read as the grade half of a morphism, specialised to the case where
one only needs the image of source canonical trust under $\psi$.

\paragraph{Scope.}
The development treats a single map of chains and two finite examples. It
does not construct a full morphism of trust systems (claim and step
components), nor does it address naturality of transport under composition
of morphisms. Those belong to the general theory.

\paragraph{Machine checking.}
Every definition and theorem statement above corresponds to a declaration in
a machine-checked formalisation. The phrase ``by lean'' indicates that the
proof has been verified by a proof assistant; we have not reproduced the
formal proof scripts here. A mathematician reading this note need only
verify that the English statements match the intended mathematics; the
formal system is responsible for the derivations.

%======================================================================
\section*{Acknowledgements}
%======================================================================

The mathematical content composes the trajectory systems of the companion
note with the chain-morphism requirements of the theory of trust systems.
The presentation here is a prose rendering of a formal development intended
for human review of the theorem statements.

\begin{thebibliography}{9}

\bibitem{trust-systems}
Trust systems: the certificate lemma and transport of trust.
\newblock Companion formal development.

\bibitem{trajectory}
Trajectory systems: circular verification collapses canonical trust.
\newblock Companion note.

\end{thebibliography}

\end{document}
